% ==========================================
% BAB IV DESAIN KONSEP SOLUSI
% ==========================================
\chapter{DESAIN KONSEP SOLUSI}
\label{chap:desain-konsep-solusi}

\section{Gambaran Umum Solusi}

Berdasarkan analisis pada Bab III, solusi yang diajukan adalah:
\begin{itemize}
    \item Implementasi infrastruktur edge computing berbasis Single Board Computer untuk hosting aplikasi ERP.
    \item Evaluasi kelayakan teknis dan finansial deployment ERP pada perangkat bertenaga rendah sebagai alternatif terhadap cloud konvensional.
    \item Framework pengujian untuk perbandingan performa dan biaya antara deployment lokal dan cloud.
\end{itemize}

\section{Arsitektur Deployment}

\subsection{Deployment pada Edge Computing}
\begin{itemize}
    \item Konsumsi daya rendah dengan investasi satu kali tanpa biaya berlangganan berkala.
    \item Kontrol penuh atas data dengan kemampuan operasi offline.
    \item Platform menggunakan sistem operasi berbasis Linux dengan stack aplikasi standar.
    \item Akses jaringan lokal untuk integrasi dengan perangkat bisnis.
\end{itemize}

\subsection{Deployment pada Cloud}
\begin{itemize}
    \item Managed infrastructure dengan skalabilitas resources sesuai kebutuhan.
    \item High availability dengan biaya operasional berulang.
    \item Dependency pada koneksi internet stabil.
    \item Digunakan sebagai baseline perbandingan performa dan biaya.
\end{itemize}

\subsection{Simulasi Distributed System}
\begin{itemize}
    \item Simulasi distributed deployment menggunakan metode alternating karena keterbatasan sumber daya.
    \item Membuktikan kelayakan multi-node deployment dengan single unit hardware.
\end{itemize}

\section{Metodologi Pengujian}

\subsection{Skenario Pengujian}
\begin{itemize}
    \item Load testing untuk mensimulasikan workload UMKM.
    \item Simulasi concurrent users dengan transaction scenarios umum.
    \item Pengukuran performa pada berbagai kondisi beban.
\end{itemize}

\subsection{Metrics Pengukuran}
\begin{itemize}
    \item Response time untuk berbagai operasi bisnis.
    \item Throughput dan kapasitas pemrosesan transaksi.
    \item Resource utilization (CPU, memori, storage).
    \item Error rate dan reliability sistem.
\end{itemize}

\subsection{Infrastruktur Monitoring}
\begin{itemize}
    \item Tools monitoring untuk data collection real-time.
    \item Analisis performa berdasarkan metrics yang dikumpulkan.
\end{itemize}

\section{Analisis Biaya}

\subsection{Total Cost of Ownership (TCO)}
\begin{itemize}
    \item Perbandingan TCO untuk deployment edge computing vs cloud.
    \item Capital expenditure: investasi awal hardware dan software.
    \item Operational expenditure: biaya listrik, maintenance, dan subscription.
\end{itemize}

\subsection{Proyeksi Finansial}
\begin{itemize}
    \item Proyeksi biaya multi-tahun untuk kedua alternatif infrastruktur.
    \item Identifikasi break-even point antara investasi edge computing dan cloud.
    \item Analisis cost-effectiveness berdasarkan skala dan kebutuhan UMKM.
\end{itemize}
